\section{Pairing two character tails before estimating them}
\label{sz-tails-section}

There is another useful simplification in Guo's paper \cite{sz-guo}.
It avoids a spurious factor $(1-\beta)^{-1}$ in a hyperbola estimate.
The next version applies to any bounded sequence with bounded interval sums
and records the best constant supplied by this particular pairing.

\begin{proposition}\label{sz-paired-tails}
Let $c:\mathbb N\to\mathbb C$ satisfy $|c(n)|\le1$ and
\[
 \left|\sum_{u<n\le v}c(n)\right|\le B\qquad(0\le u\le v).
\]
For real $s>0$, set $L_c(s)=\sum_{n\ge1}c(n)n^{-s}$, which converges
by summation by parts. Suppose $0<\beta<1$ and $L_c(\beta)=0$.
Put $\varepsilon=1-\beta$ and $a(n)=\sum_{d\mid n}c(d)$.
Then, for every $t\ge1$,
\begin{align}\label{sz-paired-main}
 \sum_{n\le t}\frac{a(n)}{n^\beta}
 &=\frac{L_c(1)}{\varepsilon}t^\varepsilon+E(t),\\
 |E(t)|&\le
 \left[B\bigl(1+2\beta^{\beta/(1-\beta)}\bigr)
                  +1+\tfrac12\log t\right]t^{1/2-\beta}.\nonumber
\end{align}
For $\beta\ge1/2$, the bracket can be replaced by
$2B+1+\tfrac12\log t$.
\end{proposition}
\begin{proof}
Write $A_u(v)=\sum_{u<n\le v}c(n)$. Summation by parts gives
\[
 R_s(u):=\sum_{n>u}\frac{c(n)}{n^s}
      =s\int_u^\infty A_u(v)v^{-s-1}\,dv,
 \qquad |R_s(u)|\le Bu^{-s}.
\]
Set $u=\sqrt t$. Splitting the pairs $dm\le t$ according to $d\le u$
or $d>u$ gives
\[
 \sum_{n\le t}\frac{a(n)}{n^\beta}
 =\sum_{d\le u}d^{-\beta}\sum_{m\le t/d}\frac{c(m)}{m^\beta}
  +\sum_{m\le u}\frac{c(m)}{m^\beta}
            \sum_{u<d\le t/m}d^{-\beta}.
\]
The first sum has modulus at most $Bt^{1/2-\beta}$, because its inner
sum is $-R_\beta(t/d)$. For real $v\ge u\ge1$, integral summation gives
\[
 \sum_{u<n\le v}n^{-\beta}
 =\frac{v^\varepsilon-u^\varepsilon}{\varepsilon}+r(u,v),
 \qquad |r(u,v)|\le u^{-\beta}.
\]
The error contributed to the second sum is therefore at most
\[
 u^{-\beta}\sum_{m\le u}m^{-\beta}
 \le u^{\varepsilon-\beta}(1+\log u).
\]
Its two main terms are
\[
 \frac{L_c(1)}\varepsilon t^\varepsilon
 -\frac{t^\varepsilon}\varepsilon R_1(u)
 +\frac{u^\varepsilon}\varepsilon R_\beta(u).
\]
Pair the last two terms before taking absolute values. Their sum is
\[
 u^\varepsilon\sum_{m>u}c(m)g(m),\qquad
 g(v)=\frac{v^{-\beta}-u^\varepsilon v^{-1}}\varepsilon.
\]
The function $g$ is nonnegative, vanishes at $u$ and at infinity, and
increases and then decreases, with maximum at $v=u\beta^{-1/\varepsilon}$.
A direct substitution gives
$\max g=u^{-\beta}\beta^{\beta/\varepsilon}$.
Consequently its total variation is twice this maximum. A final summation
by parts bounds the paired terms by
\[
 2B\beta^{\beta/\varepsilon}u^{\varepsilon-\beta}.
\]
Adding the three errors gives \eqref{sz-paired-main}. Finally
$\beta^{\beta/(1-\beta)}$ decreases from $1/2$ to $e^{-1}$ on
$[1/2,1)$, since the derivative of its logarithm is
$(\log\beta+1-\beta)/(1-\beta)^2\le0$.
\end{proof}

For a nonprincipal real character, one may take $B=q/2$ by deleting complete
periods from an interval and comparing the remaining block with its
complement. In that case $a=1*\chi$ is nonnegative: at a prime power,
\[
 a(p^k)=
 \begin{cases}
 k+1,&\chi(p)=1,\\
 1,&\chi(p)=0,\\
 (1+(-1)^k)/2,&\chi(p)=-1.
 \end{cases}
\]
This is why the following comparison applies to real characters.

\begin{lemma}[A weighted positive-sum comparison]\label{sz-positive}
Suppose that $\chi$ is a nonprincipal real character,
$1/2\le\beta<1$, and $L(\beta,\chi)=0$. Put $\varepsilon=1-\beta$,
\[
 H_1(y)=\sum_{n\le y}\frac{(1*\chi)(n)}n,\quad
 B_0(t)=(q+1+\tfrac12\log t)t^{-1/2},\quad
 K=\log x\log y+\tfrac12(\log x)^2.
\]
If $1\le y\le x$ and $\varepsilon\log(xy)\le1$, then
\begin{align*}
 \sum_{y<p\le x}\frac{(1+\chi(p))\log p}{p}
 \le\varepsilon K+\frac{1}{H_1(y)}
 \bigl[B_0(y)(\varepsilon K+\log y)
       +B_0(xy)\log(xy)+2B_0(y)+2y^{-1/2}\bigr].
\end{align*}
\end{lemma}
\begin{proof}
Write $a=1*\chi$ and $N(t)=\sum_{n\le t}a(n)(\log n)/n$.
The map $(n,p)\mapsto np$, with $n\le y<p\le x$, is injective, since
equality $np=n'p'$ with distinct primes would force $p\mid n'$, although
$n'<p$. Multiplicativity and nonnegativity give
\[
 H_1(y)\sum_{y<p\le x}\frac{a(p)\log p}{p}\le N(xy)-N(y).
\]
Use Proposition~\ref{sz-paired-tails} with $c=\chi$ and put
$D(t)=t^{-\varepsilon}E(t)$; thus $|D(t)|\le B_0(t)$.
Partial summation, applied to $(\log t)t^{-\varepsilon}$, gives exactly
\begin{align*}
 N(xy)-N(y)&=L(1,\chi)K+R,\\
 R&=D(xy)\log(xy)-D(y)\log y
      +\int_y^{xy}(\varepsilon\log t-1)D(t)\frac{dt}{t}.
\end{align*}
Because $0\le\varepsilon\log t\le1$ on the interval,
\[
 |R|\le B_0(xy)\log(xy)+B_0(y)\log y
             +\int_y^\infty B_0(t)\frac{dt}{t}.
\]
The last integral is $2B_0(y)+2y^{-1/2}$. Also, positivity gives
\[
 \frac{L(1,\chi)}\varepsilon
 =y^{-\varepsilon}\left(\sum_{n\le y}\frac{a(n)}{n^\beta}-E(y)\right)
 \le H_1(y)+B_0(y).
\]
Substitute this inequality and divide by $H_1(y)\ge a(1)=1$.
\end{proof}

The pairing is useful whenever a proof needs a uniform error as a real zero
approaches $1$. It does not produce high multiplicity for an automorphic
theta lift, nor does positivity of $1*\chi$ extend to arbitrary complex
characters. Those are distinct issues in a proof of the full generalized
Riemann hypothesis.

